\section{Lean implementation plan and concrete dependencies}
\label{sec:integration-new}

\subsection{Use the existing outer boundary}
Forge already recommends starting with ordinary Lean goals and feeding checked
facts to a fresh local \code{grind} run. The current Lean documentation describes
\code{grind} as cooperating local solvers with proof construction; these workers
should augment that machinery, not reimplement it~\cite{forge-integration,grind}.
No new shared congruence-closure database is required for the first version.

Use four focused modules behind the existing worker interface:
\begin{lstlisting}
Forge.Worker.InvariantSpace    // constant and polynomial multipliers
Forge.Worker.FiniteCover       // cover or concrete tree
Forge.Worker.IntProjection     // guard and uniform witness
Forge.Planner.ContractSketch  // proof-directed hole decomposition
\end{lstlisting}
These are proposed module names. None is present as an implemented Lean worker
in this archive.

\subsection{A minimal certificate-to-proof path for each worker}
\paragraph{Invariant packets.}
Reification returns a polynomial syntax for each source expression together with
a theorem relating its evaluation to that expression. A sparse normalizer proves
that equality of canonical coefficient maps implies equality of evaluations.
The checker validates initial substitutions and multiplier identities, then
instantiates a generic reachability-induction theorem. To derive a particular
goal equality from the packet, supply and check its coordinates as a constant
or polynomial combination of the generators. A packet is useful only after this
last target connection is made.

For rational coefficients interpreted in an ordered field or another
characteristic-zero field, use an explicit coefficient homomorphism. For integer
source states, clearing denominators is often simpler than transporting many
rational divisions. Neither method licenses dropping a source cast or treating
natural subtraction as integer subtraction. The reifier must reject unsupported
operations rather than silently interpreting them as polynomial operations.

\paragraph{Finite covers.}
A positive certificate is a finite set of states; a reflected checker establishes
the constructor closure and good-state inclusion. A generic structural theorem
then yields the source property using the registered summary laws. A negative
certificate is a constructor DAG, evaluated and reconstructed as a concrete
source term. The two result forms must have different proof goals.

The included \code{CoreSoundness.lean} states the generic reachability and tree
cover theorems in a small core-only form. It does not include the reflected
finite checker or source reifier. Finite verification can begin with ordinary
\code{decide} on small explicit instances; larger arrays may justify a dedicated
reflected checker with a proved interpreter. An unchecked native computation
must not be substituted for that theorem.

\paragraph{Integer projection.}
The first implementation can elaborate each receipt by theorem application rather
than formalizing the entire search algorithm. Prove the congruence normalization,
CRT merge, floor-bound, and nearest-point lemmas once. Decode the program as a
shared expression with positive constant divisors. Return both an equivalence
proof for the feasibility guard and a proof for the selected witness under that
guard. Keeping only the second loses useful negative and routing information;
keeping only the first loses constructive synthesis.

\paragraph{Contract sketches.}
A registered schema supplies the proof of composition. Hole candidates are
elaborated at exact source types and checked against their full contracts.
Assemble the final expression and theorem only after all contracts are proved.
The source-typed Djex boundary is a sensible candidate-provider API; the
behavioral observations in its diagnostic reports are filters, not replacements
for these universal contracts~\cite{djex-readme,leant-readme}.

\subsection{Library choices, with their roles kept narrow}
\begin{center}
\small
\begin{tabularx}{\textwidth}{@{}p{0.22\textwidth}YY@{}}
\toprule
Component & Concrete dependency & Role and boundary \\
\midrule
Local proof closure & Lean \code{grind}, \code{omega}; Mathlib \code{ring} & Close ordinary proof obligations after explicit side conditions are supplied \\
Polynomial semantics & Mathlib \code{MvPolynomial}, especially evaluation homomorphisms & Source denotation and correctness lemmas; not external search authority \\
Finite certificates & Lean/Mathlib \code{Array}, \code{Fin}, \code{Fintype}, \code{Finset} & Enumerate finite states and prove a checker correct \\
Exact proposer algebra & \code{fractions.Fraction} now; optional \code{python-flint.fmpq_mat} & Nullspaces and rational solving outside the kernel \\
Integer proposer & Extended Euclid plus the included DAG assembler & Deterministic construction; every arithmetic lemma reconstructed in Lean \\
Typed component search & Djex's checked public synthesis boundary & Supply source-owned hole candidates, preserving polymorphism and providers \\
Optional grammar search & cvc5 SyGuS & Propose first-order arithmetic candidates; replay their contracts in Lean \\
\bottomrule
\end{tabularx}
\end{center}

The table names implementation components, not a tested dependency matrix.
Mathlib's current documentation uses the path
\path{Mathlib/Algebra/MvPolynomial/Eval.lean}; names and imports must be checked
against the project's pinned Mathlib revision, not inferred from a historical
\code{Data/MvPolynomial} path~\cite{mvpoly}. The pinned Forge baseline reports
Lean 4.34.0 and a specific Mathlib commit; the Python prototypes do not depend on
that installation and do not establish adapter compatibility with it.

\subsection{Mathematical contracts for the native interface}
The prepared problem should describe which interpretation theorem the worker is
allowed to produce. A useful separation is:
\begin{lstlisting}
InvariantProblem -> packet identities -> reachable-state theorem
FiniteProblem    -> cover proof | concrete-tree proof
IntegerProblem   -> (guard equivalence, witness correctness)
HoleProblem      -> candidate term with full contract proof
\end{lstlisting}
A closed packet about an auxiliary transition system is not yet a proof of an
unrelated target. A finite cover for an abstraction is not a concrete
counterexample. A projection guard is not an existential witness until its
premise is proved. These are worker-specific proof obligations beyond the
identity and scope discipline already specified by Forge.

For replay, retain the re-checkable problem and certificate, not a producer's
cached verdict. The existing Forge contracts for exact source identity,
rollback, and proof acceptance remain applicable; this article does not replace
them with a second protocol.

\subsection{One particularly useful proof-reuse boundary}
Compile a successful search into a theorem specialized to its transition maps,
finite algebra, or arithmetic coefficients. The general checker theorem can be
shared, while the certificate itself becomes finite checked data. Recompilation
then need not rerun Gaussian elimination, reachability search, or a synthesis
solver. This applies equally to a reusable hole theorem from the sketch planner.

Avoid serializing expanded polynomial normal forms into a huge source script
when a coefficient vector and a generic normalization proof suffice. Conversely,
for a small example, a short ordinary proof may be simpler and faster than a
new reflected interpreter. Measure both source size and kernel-checking cost at
the actual target before choosing the permanent representation.

\section{Executed experiments and what they establish}
\label{sec:experiments}

\subsection{Environment and units}
The recorded run used Python 3.13.5 on Linux x86-64. The command line used
\code{python -S}, disabling site-package initialization. The implementations
import only the standard library. Neither \code{lean} nor \code{lake} was
available, and the Lean status file records \textsc{not run} rather than a
successful elaboration.

There is one new experiment ledger, separate from Forge's nine historical runs.
A certificate, a unit-test method, a parameter valuation, a state, and a tuple
evaluation are different units. The tables keep them separate. In particular,
these records must not be appended to a historical solved-goal total.

\begin{center}
\small
\begin{tabularx}{\textwidth}{@{}lrrY@{}}
\toprule
Lane & Stored certificates & Replay accepted & Contents \\
\midrule
Constant multipliers & \ConstCertificates & \ConstCertificates & Cyclic packet, recurrence, 48 affine cases, one expected miss \\
Polynomial multipliers & \IdealCertificates & \IdealCertificates & Squaring, union, parabola, false-line control, powers 2 through 9 \\
Finite algebra & \FiniteCertificates & \FiniteCertificates & Count covers, mirror examples, and random exact algebras \\
Integer projection & \IntegerCertificates & \IntegerCertificates & 220 random symbolic systems and one large-period system \\
\bottomrule
\end{tabularx}
\end{center}

There are 422 stored certificates in this one corpus. Some are concrete
counterexamples and some contain an empty invariant packet. The count is not
422 proved Lean theorems, and it is not a solved-goal benchmark.

\subsection{Invariant tests and ablations}
The cyclic-differences case returns a dimension-two invariant space while its
conserved-linear subspace has dimension zero. The nonlinear squaring-diagonal
case returns no generator with multiplier degree zero and one generator with
multiplier degree one. The union-of-axes case changes from zero to one generator
when the multiplier bound rises from one to two. The false-line control confirms
that an initially true relation is not retained merely because it fits the
initial parameterization.

The 48 affine-diagonal instances have three state variables, one to three
transitions, and alternating template degrees one and two. Random affine maps
are constructed to preserve the diagonal. The computed dimensions are two and
seven, respectively. These engineered cases exercise simultaneous invariants,
multiple transitions, and coefficient handling; they are not a representative
sample of all programs.

\begin{center}
\small
\begin{tabularx}{\textwidth}{@{}YrrY@{}}
\toprule
Example & $d$ & $e$ & Recorded result \\
\midrule
Cyclic differences & 1 & 0 & Two generators; no nonzero conserved linear generator \\
Sums of squares & 3 & 0 & One generator after dimensions $9,5,3,2,1,1$ \\
Squaring, initial $x=y$ & 1 & 0 & Empty packet: expected certificate-language miss \\
Squaring, initial $x=y$ & 1 & 1 & $y-x$ with multiplier $x+y$ \\
Squaring, initial $xy=0$ & 2 & 1 & Empty packet: multiplier bound too small \\
Squaring, initial $xy=0$ & 2 & 2 & $xy$ with multiplier $xy$ \\
Squaring, initial $y=x^2$ & 2 & 2 & $y-x^2$ with multiplier $y+x^2$ \\
\bottomrule
\end{tabularx}
\end{center}

Displayed generators can differ by sign or rational scaling from the stored
row-reduced basis. The actual basis and exact multiplier matrices remain in the
certificate file. The table explains their mathematics rather than rewriting
the recorded evidence.

\subsection{Finite-algebra differential checks}
The count family covers moduli 1 through 12. Each case returns a cover with
exactly $m$ states in a product of size $m^2$. The mirror family has 24 magma
instances: 12 commutative cases return covers and 12 deliberately noncommutative
cases return counterexample trees. The 100 random finite algebras return 33
covers and 67 counterexamples, agreeing with a separately written full-closure
oracle.

The oracle computes all reachable states even when the production proposer has
already found a bad state. Thus the comparison checks the verdict against the
whole exact algebra, not only against the first explored tree. The positive
certificate checker separately enumerates all constructor tuples in the
returned cover. The negative checker independently evaluates its acyclic tree.
Both still share input-table semantics with the search, a much smaller shared
surface than the worklist algorithm but not a formal verification of it.

\subsection{Integer differential checks}
The 220 random symbolic systems use one or two parameters, fixed output
coefficients in $[-4,4]$, moduli from 1 through 8, up to four inequalities, and
up to three congruences. At each parameter valuation from
$\{-4,-1,0,2,5\}^n$, the generated feasibility guard is compared with a complete
finite enumeration for that \emph{concrete} problem.

Let $B$ be one plus the largest absolute right-side value, and let $P$ be the
lcm of the input moduli, or one if there are none. Every finite bound endpoint
lies in $[-B,B]$. The congruence solution set is periodic with period dividing
$P$. Therefore any satisfiable interval, including a one-sided or fully
unbounded interval, has a witness in $[-B-P,B+P]$. Enumerating that range is a
complete concrete oracle for the generated test inputs, not merely an arbitrary
small box.

Across 3,400 valuations, the guard reports 1,368 feasible and 2,032 infeasible
instances, all agreeing with the oracle. Every feasible valuation also checks
the selected witness directly against the original input constraints. The
large-period case is tested separately at five values, including
$\pm10^{100}$; it is not subjected to impossible modulus-sized enumeration.

\subsection{Corruption and implementation controls}
All 21 unit-test methods pass. Their internal cases include malformed rational
coefficients, matrix and multiplier corruption, wrong initial behavior, changed
problem subjects, an insufficient multiplier budget, an empty or unsafe cover,
a self-referential tree, a good state falsely designated as a counterexample,
changed Bezout coefficients, deleted feasibility guards, a changed witness node,
and floating-point or Boolean values where exact integers are required.
There are 19 explicit certificate-mutation cases in those methods, in addition
to malformed-polynomial and semantic differential cases.

The replay script disables synthesis, row reduction, nullspace, coordinate
search, and extended-gcd construction by replacing their entry points with
exceptions before checking the stored corpus. Every certificate still replays.
This demonstrates separation from those search routines; it does not make
shared polynomial arithmetic or the integer assembler formally verified.

\subsection{Timings are diagnostic only}
The raw ledger includes one observed duration per lane. In the recorded run,
constant-multiplier experiments took approximately 0.613 seconds, polynomial
multipliers 0.037 seconds, finite algebras 0.023 seconds, and integer projections
0.260 seconds. These figures include that lane's construction, checks, and
specified oracles, but exclude process startup and final JSON serialization.
The lanes perform different work, and there are no repeated trials, confidence
intervals, or matched Lean baselines. The figures are useful for detecting gross
reproduction problems, not for claiming comparative performance.

\section{Implementation order and acceptance gates}
\label{sec:roadmap}

\subsection{First: one source theorem per certificate family}
Begin by elaborating the supplied generic Lean specimens under the pinned
project and implementing one actual source bridge per family. A good minimum
set is the cyclic packet, a binary-tree count congruence, and a guarded
non-coprime integer witness. Require a complete proof of the exact source goal,
with recorded axiom dependencies, and then rerun it without the proposer.
This establishes something the Python corpus cannot: a certificate really can
cross the source-language boundary into a Lean proof.

\subsection{Second: native certificate checkers, not native searches}
Port the deterministic checker first. The external Python proposer can continue
to supply exact finite data. For finite covers, prove a table-closure checker
sound. For integer projection, compose the four generic arithmetic lemmas. For
invariant packets, finish polynomial reification and identity checking. Only
then optimize or reimplement the external searches in Lean, FLINT, or another
language if measurements justify it.

The polynomial multiplier algorithm is mathematically the most substantial
addition, but its source polynomial bridge is also the most delicate. The finite
cover checker is a useful smaller first end-to-end milestone. This is an
engineering order, not a ranking of the algorithms' eventual importance.

\subsection{Third: benchmark the actual new capability}
Build source-goal families that isolate the extension: simultaneous invariants
without conserved generators; invariant ideals requiring positive multiplier
degree; arbitrarily deep trees summarized by a finite algebra; and integer
witnesses with large symbolic periods. Include true goals outside the supported
fragments and false goals whose counterexamples are checked.

Compare against \code{grind} with the same definitions and permitted theorem
inventory, a reasonable hand-orchestrated baseline, and relevant specialized
tactics. Report separately preparation, search, certificate size, reconstruction,
kernel checking, and cached replay. Hold out theorem families rather than only
renaming variables, and prevent retrieval from importing the exact target
lemma. None of those controlled Lean measurements is claimed by this article.

\subsection{Fourth: test the original higher-order query unchanged}
Integrate the contract-sketch branch into the actual Leant/Djex path and run the
original indexing query with its original source type, provider inventory,
integer semantics, observations, and resource limits. Retain exact branch traces
showing carrier choice, hole inventories, and candidate order. The list-fold
specimen and isolated step successes are useful controls; neither substitutes
for this gate. The simultaneous two-universe query remains a separate target.

\section{Assessment}
The strongest addition is the bounded-multiplier fixed point. It provides a
finite, exact way to synthesize entire invariant packets, rather than asking a
bilinear solver to guess generators and their closure multipliers at once.
Finite-algebra covers give another precise form of induction: a small closure
certificate stands for all finite constructor trees. Constructive projection
turns one-output integer feasibility into a compact program instead of a residue
table. Contract-first sketches explain how these mathematical workers can
cooperate with a sophisticated type-directed synthesizer without reducing the
problem to a larger unstructured candidate stream.

The mathematical claims are deliberately local and useful. The invariant worker
is complete for its bounded certificate space; affine transitions admit the
stronger degree-bounded equality result. The finite worker decides its explicit
finite algebra. The integer worker exactly projects its stated one-output
conjunctive fragment. None is a complete general Lean prover, and none needs to
be one to add a qualitatively different successful proof route.

The executable archive establishes that the core constructions work on the
recorded cases, reject the specified corruptions, and replay without their
search routines. The next decisive result is not a larger Python pass count: it
is a checked source theorem from each family, followed by a fair measurement in
the actual Forge environment. The article supplies the algorithms, certificate
contracts, proofs, concrete implementation choices, and reproducible artifacts
needed to take that step.
